% Chapter 0: 综合评述（批判性分析）
% 路径：/Users/yueyh/Projects/aca-workflow/notes/information-geometry/tang-rong/chapters/chapter0.tex
\documentclass[12pt]{amsbook}
\usepackage{amsmath,amssymb,amsthm}
\usepackage{xeCJK}
\usepackage{microtype}
\usepackage{hyperref}
\usepackage{cleveref}
\usepackage{geometry}
\usepackage{natbib}
\bibliographystyle{plainnat}
\geometry{margin=1in}

\theoremstyle{plain}
\newtheorem{Definition}{定义}[chapter]
\newtheorem{Theorem}{定理}[chapter]
\newtheorem{Lemma}{引理}[chapter]
\newtheorem{Remark}{注}[chapter]

% noteinfo: inline remark box
\usepackage{mdframed}
\usepackage{tcolorbox}
\newtcolorbox{noteinfo}{colback=yellow!15, colframe=orange!50,
  boxrule=0.3pt, left=8pt, right=8pt, top=6pt, bottom=6pt}
\NewDocumentCommand{\noteinfo}{m}{\begin{noteinfo}#1\end{noteinfo}}

% ----------------------------------------------------------------
% Standalone guard: 主笔记包含本章时跳过 document 环境
\ifdefined\STANDALONE
\else
\begin{document}
\fi
% ----------------------------------------------------------------

\chapter{综合评述}\label{ch:chapter0}

本章节综述 Rong Tang 与 Yun Yang 的系列工作，
涵盖 2019--2025 年间发表在 Annals of Statistics、JRSS B、
AISTATS、ICML 等顶级期刊与会议上的代表性论文，
涉及极小极大理论、稳健贝叶斯、扩散模型、两样本检验、MCMC 等多个方向。

\section{研究背景与动机}\label{sec:ch0-background}

现代高维数据（图像、文本）尽管环境维度 $D$ 巨大，
但其有效信息维度 $d \ll D$，数据通常位于某个低维流形上。
传统非参数统计受维度灾难困扰——极小极大速率以 $D$ 为指数，
而 GAN、VAE、扩散模型等深度生成模型在实践中却表现出色。
这一矛盾促使一个根本性问题：

\begin{enumerate}
    \item 当数据位于未知低维流形上时，\textbf{分布估计的统计极限是什么？}
    \item \textbf{深度生成模型能否隐式自适应流形结构，避免维度灾难？}
    \item \textbf{如何在流形约束下进行有效的贝叶斯推断和 MCMC 采样？}
\end{enumerate}

Tang-Yang 系列围绕这三个问题，在统计学习理论与方法之间建立了
有机联系：从极小极大理论出发，理解扩散模型的自适应性质，
并在此基础上建立稳健贝叶斯推断和高效采样算法。

\noteinfo{本章侧重批判性分析：技术创新来源、
证明策略选择、与现有工作的对话、以及开放问题。
与第 \ref{ch:tang-rong-review} 章的系统综述形成互补。}

\section{论文系列详解}\label{sec:ch0-papers}

% ----------------------------------------------------------
\subsection{\citep{tang2022minimax} --- 未知子流形上的极小极大分布估计}
\label{sec:ch0-annals2022}

\textbf{文献信息}：Rong Tang and Yun Yang.
``Minimax Rate of Distribution Estimation on Unknown Submanifold under
Adversarial Losses.'' \textit{Annals of Statistics}, 2022.

\subsubsection{Summary}

本文研究了高维数据低维流形结构假设下的分布估计问题，建立了对抗损失
（adversarial loss）下未知子流形上分布估计的极小极大（minimax）
理论框架。核心贡献在于证明了估计误差率为
\[
n^{-1/2} \;\vee\; n^{-\frac{\alpha+\gamma}{2\alpha+d}}
\;\vee\; n^{-\frac{\gamma\beta}{d}},
\]
其中 $d$ 为流形内在维度，$\alpha$ 为密度光滑度，
$\beta$ 为流形光滑度，$\gamma$ 为判别器光滑度。
该速率仅依赖内在维度 $d$ 而非环境维度 $D$，成功规避了维度灾难。

技术路线方面：下界证明依赖于精细的 Fano 方法和 Le Cam 方法，
通过在球面上构造难以区分的分布族完成；
上界的证明综合运用了小波展开、经验过程理论（chaining 技术）、
单位分解和最优传输的 Caffarelli 正则性理论。

\subsubsection{Overall Assessment}

\textbf{理论贡献}：本文成功填补了生成模型分布估计与流形统计推断
之间的理论空白。核心发现是三类不同的收敛机制（参数机制、
密度估计机制、流形估计机制）的识别，揭示了判别器光滑度 $\gamma$
在其中的关键调控作用。

\textbf{批评}：
（1）下界构造中基于对 $d$ 维球面的特定扰动是否真正代表了所考虑
分布类中的最难情形，尚有进一步探讨的空间。
（2）估计器的理论构造涉及小波展开截断、Lepski 自适应等复杂步骤，
未给出具体算法实现——这是一篇"理论导向"而非"方法导向"的工作。
（3）主要处理无噪声情形，对带噪声观测的讨论不够系统。

\textbf{在系列中的位置}：
本文是整个系列的奠基之作，后续 \citep{tang2023lowdim}、
\citep{tang2024diffusion}、\citep{tang2025regression}
的工作都在此极小极大框架下被评价和比较。

% ----------------------------------------------------------
\subsection{\citep{tang2021vae} --- VAE 的 Excess Risk Bound}
\label{sec:ch0-colt2021}

\textbf{文献信息}：Rong Tang and Yun Yang.
``Excess Risk Bounds for Distribution Estimation with VAE.'' \textit{COLT}, 2021.

\textbf{核心发现}：VAE 的 encoder-decoder 结构本质上是在学习
数据流形的一个全局参数化。建立了 VAE 估计器的 excess risk 上界，
证明了在温和假设下 VAE 可以达到最优的统计收敛速率。

\textbf{批评}：
（1）假设条件依赖于解码器族和编码器族的逼近能力，
在实际应用中的可验证性有待进一步讨论。
（2）VAE 的 posterior collapse 问题在理论上没有得到解释，
限制了理论结果的实践指导意义。
（3）作为系列的开创性工作，理论框架相对简单，
主要结果是 excess risk bound 而非极小极大速率比较。

% ----------------------------------------------------------
\subsection{\citep{tang2023robust} --- 黎曼子流形上的稳健贝叶斯推断}
\label{sec:ch0-jrssb2023}

\textbf{文献信息}：Rong Tang, Anirban Bhattacharya, Debdeep Pati,
and Yun Yang. ``Robust Bayesian Inference on Riemannian Submanifold,''
\textit{JRSS B}, 2023.

\textbf{核心理论贡献：黎曼流形上的 Bernstein-von Mises (BvM) 定理}。
通过将后验分布投影到流形切空间并利用局部坐标变换，
成功将欧氏空间中的分析工具推广到流形 setting，
证明了后验分布在切空间上的投影渐近正态。

\textbf{计算贡献：黎曼随机游走 Metropolis (RRWM) 算法}。
证明了混合时间仅依赖流形内在维度 $d$ 而非环境维度 $D$——
分析工具是经典的 conductance profile 方法。

\textbf{批评}：
（1）正则性假设较强：假设流形局部为 $C^3$ 光滑，
且要求风险函数满足二次增长条件，限制了在更一般流形上的应用。
（2）BvM 定理的收敛速度（到渐近正态的速率）在论文中没有给出，
使得有限样本下的可靠性难以量化。
（3）RRWM 的 proposal 设计较为简单，在多峰分布或强非线性流形上
可能表现不佳。

% ----------------------------------------------------------
\subsection{\citep{tang2024mala} --- Metropolis-Adjusted Langevin 算法的最优混合时间}
\label{sec:ch0-icml2024}

\textbf{文献信息}：Rong Tang and Yun Yang.
``Metropolis-Adjusted Langevin Algorithm: Optimal Dimension Dependency
and Beyond.'' \textit{ICML}, 2024.

\textbf{核心贡献：$s$-conductance profile}。
这是介于经典 conductance 与 Chen et al. (2020) 的 conductance profile
之间的精细分析技术。核心创新：在非凸集假设下
将 warm-start 参数依赖从 $\log M_0$ 改进到 $\log\log M_0$——
在高维后验分布（$M_0$ 极大）的实际场景中意义显著。

在无需对目标分布施加光滑或强对数凹性假设的条件下，
证明了 MALA 达到最优 $O(d^{1/3})$ 维度依赖（经 burn-in 后），
并将结果推广到光滑与非光滑 Gibbs 后验（包括贝叶斯分位数回归）。

\textbf{批评}：
（1）定理要求 $d \leq c n^{\kappa_1}/\log n$，对于非光滑情形
需要相对更大的样本量才能保证 $d^{1/3}$ 混合时间。
（2）$s$-warm start 假设在实践中难以验证，
与 MCMC "无知采样"的哲学存在张力。
（3）没有给出实际可用的 burn-in 长度估计。

\textbf{与 \citep{tang2023robust} 的联系}：
\citep{tang2023robust} 的 BvM 定理保证了流形后验的渐近高斯性；
本文则利用这一性质，将 MALA 在高斯目标上的最优混合时间
桥接到一大类后验分布。

% ----------------------------------------------------------
\subsection{\citep{tang2023twosample} --- 对抗损失下的两样本检验}
\label{sec:ch0-aistats2023}

\textbf{文献信息}：Rong Tang and Yun Yang.
``Minimax Nonparametric Two-Sample Test under Adversarial Losses,''
\textit{AISTATS}, 2023.

\textbf{核心创新}：
建立了负 Besov 范数与对抗损失之间的等价关系（引理 1），
基于截断小波展开构造检验统计量，
计算复杂度 $O(n\log n + n^{2d/(4\alpha+d)})$
远优于 MMD 的 $O(n^2)$。

\textbf{批评}：
（1）正则性假设过强：假设密度函数具有紧支撑且一致有界，
在许多实际应用中可能难以满足。
（2）最优截断水平 $J$ 的选择依赖于未知的光滑度参数 $\alpha$——
如何构造对 $\alpha$ 自适应的检验仍是开放问题。
（3）Bootstrap 阈值的理论支撑不足。

% ----------------------------------------------------------
\subsection{\citep{tang2024diffusion} --- 扩散模型自适应流形结构}
\label{sec:ch0-aistats2024}

\textbf{文献信息}：Rong Tang and Yun Yang.
``Adaptivity of Diffusion Models to Manifold Structures,''
\textit{AISTATS}, 2024.

\textbf{Langevin 扩散分支}：
通过 Gaussian 平滑 $p_\sigma(x) = \int p_{\mathrm{data}}(y)\phi_\sigma(x-y)\,\mathrm{d}y$
使密度函数在任意时刻 $t>0$ 关于 Lebesgue 测度绝对连续。
收敛速率为 $\tilde{O}(\sigma)$，仅需分数函数的四阶矩误差——
较现有工作的 $L_\infty$ 或矩母函数假设更弱。

\textbf{正向-反向扩散分支}（主要贡献）：
设计了时间分段常数的神经网络架构来逼近时变分数函数。
在 1-Wasserstein 距离下达极小极大速率
\[
\frac{1}{\sqrt{n}} \;\vee\;
n^{-\frac{\alpha+1}{2\alpha+d}}.
\]

\textbf{批评}：
（1）Langevin 扩散分支的收敛速率弱于正向-反向扩散，
存在若干技术原因导致次优。
（2）各向同性高斯噪声可能"稀释"流形结构——
论文提及但未给出严格的量化分析。
（3）时间分段架构的分段数 $J$ 依赖于未知的光滑度参数 $\alpha$，
实践中难以选择。

\textbf{与 \citep{tang2022minimax} 的关系}：
本文证明了扩散模型能够达到 \citep{tang2022minimax} 建立的极小极大速率
（在 $\gamma=1$ 的 Wasserstein 度量下），建立了
"实践有效的工具 $\Leftrightarrow$ 理论最优的方法"这一对应关系。

% ----------------------------------------------------------
\subsection{\citep{tang2023lowdim} --- 低维结构混合生成模型}
\label{sec:ch0-icml2023}

\textbf{文献信息}：Rong Tang and Yun Yang.
``Estimating Distributions with Low-dimensional Structures Using Mixtures
of Generative Models.'' \textit{ICML}, 2023.

\textbf{核心发现：多编码器/解码器结构在流形分布估计中是必要的。}
提出了 MLDMAE（Mixture of Latent Distribution Matched Auto-Encoders）方法：
$K$ 个局部编码器-解码器对分别学习流形局部区域的参数化，
权重函数 $\rho_k$ 通过 softmax 实现光滑拼接。
在 1-Wasserstein 距离下达极小极大最优（至多对数因子）。

\textbf{批评}：
（1）聚类数 $K$ 的选择缺乏自适应理论指导，
仅通过实验说明 $K\approx 10$ 对 MNIST 有效。
（2）当 $\alpha > \beta-1$ 时，最优速率是否仍然保持极小极大最优，
论文仅提及但未严格证明。
（3）局部块之间的重叠区域如何处理在理论上没有得到充分分析。

% ----------------------------------------------------------
\subsection{\citep{tang2025regression} --- 分布回归的极小极大最优速率}
\label{sec:ch0-arxiv2025a}

\textbf{文献信息}：Rong Tang and Yun Yang.
``Minimax Optimal Rates for Regression on Manifolds and Distributions,''
arXiv, 2025.

\textbf{核心贡献}：
在三类层次递进的设定下系统建立了极小极大下界，
并提出融合对抗学习（IPM 度量）与同步最小二乘的混合估计量。
$\gamma$-Hölder IPM $d_\gamma$ 统一了全变差距离（$\gamma=0$）
和 1-Wasserstein 距离（$\gamma=1$）。

\textbf{批评}：
三层 Regime 的下界证明较为复杂，
但上界中的 $\frac{\alpha_Y}{\alpha_X}\cdot d_X$ 项的出现略显 ad hoc——
缺乏几何直觉。
此外，文中提出的混合估计量在 $\gamma \neq 0,1$ 时
是否仍然可计算（如何实现 IPM 的判别器优化？）并未讨论。

% ----------------------------------------------------------
\subsection{\citep{tang2025conditional} --- 条件扩散模型的极小极大最优性}
\label{sec:ch0-arxiv2025b}

\textbf{文献信息}：Rong Tang and Yun Yang.
``Conditional Diffusion Models are Minimax-Optimal and Manifold-Adaptive,''
arXiv, 2025.

\textbf{核心发现}：
首次建立了条件扩散模型的流形自适应理论，
在经典密度回归和高维回归（协变量与响应均具有低维流形结构）
两种场景下均达到极小极大最优收敛速率。
估计误差仅依赖内在维度 $(d_X,d_Y)$ 而非环境维度 $(D_X,D_Y)$。

\textbf{批评}：
（1）与 \citep{tang2025regression} 相比，
本文缺乏对条件生成模型与条件分布回归理论之间关系的讨论。
（2）条件扩散模型中的 classifier-free guidance 等实践技术
的理论解释尚未涉及。
（3）证明中的经验过程局部化与剥离技术虽然标准，
但技术细节高度复杂，读者友好性有待提高。

\section{批判性总结}\label{sec:ch0-critique}

\subsection{系列的整体贡献}

Tang-Yang 系列的核心贡献可以概括为三点：

\begin{enumerate}
    \item \textbf{几何 $\to$ 统计的量化桥梁}：
          首次将"数据位于低维流形"这一几何先验
          转化为可计算的极小极大速率，
          其中环境维度 $D$ 完全不出现在速率指数中。
    \item \textbf{理论工具的创造性应用}：
          单位分解、Girsanov 定理、$s$-conductance profile 等工具
          的选择具有高度针对性，每一种工具都解决了
          流形统计学中的特定困难。
    \item \textbf{理论与实践的紧密互动}：
          从极小极大理论 $\to$ VAE $\to$ 扩散模型 $\to$ 条件扩散模型，
          理论每前进一步，都伴随着具体方法的验证或改进。
\end{enumerate}

\noteinfo{该系列的整体启示：低维流形几何结构是理解现代高维数据统计推断的
关键钥匙。极小极大理论告诉我们"理论上能达到多好的结果"，
稳健贝叶斯告诉我们"如何在模型不确定时量化不确定性"，
扩散模型则告诉我们"实践中如何高效地实现这一最优性"。}

\subsection{有待解决的问题}

\begin{enumerate}
    \item \textbf{流形维数自适应}：整个系列假设内在维度 $d$ 已知。
          如何在 $d$ 未知的情况下构造自适应估计器，
          同时保持极小极小最优性？
    \item \textbf{噪声鲁棒性}：
          \citep{tang2022minimax} 的无噪声设定过于理想。
          在随机噪声（deconvolution）设定下，
          流形结构能否依然帮助克服维度灾难？
    \item \textbf{光滑性自适应}：
          小波/局部多项式逼近中的截断参数依赖于未知光滑度。
          是否存在计算可行（而非仅理论存在）的自适应方法？
    \item \textbf{深度模型的端到端理论}：
          现有理论建立在"估计器 = 理论构造"的框架上，
          与实际训练（梯度下降、随机优化）的动态过程
          存在显著差距。
    \item \textbf{非光滑流形}：
          整个系列假设流形 $\beta$-光滑（$\beta \geq 1$）。
          现实中大量数据（分片光滑流形、带边流形）尚无理论保障。
\end{enumerate}

\section{发表时间线}\label{sec:ch0-timeline}

\begin{enumerate}
    \item \textbf{2021} --- \citep{tang2021vae}（VAE 流形自适应）
    \item \textbf{2022} --- \citep{tang2022minimax}
          （流形分布估计极小极大框架）
    \item \textbf{2023a} --- \citep{tang2023lowdim}（MLDMAE）
    \item \textbf{2023b} --- \citep{tang2023robust}
          （流形稳健贝叶斯 + BvM）
    \item \textbf{2023c} --- \citep{tang2023twosample}（两样本检验）
    \item \textbf{2024a} --- \citep{tang2024diffusion}
          （扩散模型自适应流形）
    \item \textbf{2024b} --- \citep{tang2024mala}（MALA 最优混合时间）
    \item \textbf{2025a} --- \citep{tang2025regression}
          （分布回归极小极小速率）
    \item \textbf{2025b} --- \citep{tang2025conditional}
          （条件扩散模型最优性）
\end{enumerate}

\noteinfo{
\textbf{主线一}：\citep{tang2022minimax} 极小极大理论
$\;\xrightarrow{\;\text{具体化}\;}\;
\citep{tang2023lowdim} \xrightarrow{\;\text{推广}\;}
\citep{tang2025regression}

\textbf{主线二}：\citep{tang2021vae}
$\;\xrightarrow{\;\text{扩展到流形}\;}\;
\citep{tang2024diffusion}
$\;\xrightarrow{\;\text{推广到条件设定}\;}\;
\citep{tang2025conditional}

\textbf{主线三}：\citep{tang2023robust}
$\;\xrightarrow{\;\text{利用 BvM 性质}\;}\;
\citep{tang2024mala}
}

\section{核心技术索引}\label{sec:ch0-tech-index}

\begin{enumerate}
    \item \textbf{单位分解（Partition of Unity）}：
          \citep{tang2022minimax}、\citep{tang2023lowdim}
          利用此技术处理无全局参数化流形——
          是估计器构造的核心工具。
    \item \textbf{$\gamma$-Hölder IPM}：
          \citep{tang2022minimax}、\citep{tang2023twosample}、
          \citep{tang2025regression} 将不同光滑度的判别器类
          统一在同一个框架下，实现多种距离度量的同时最优。
    \item \textbf{Gaussian 平滑 + 四阶矩误差}：
          \citep{tang2024diffusion} 将 Langevin 扩散的分析弱化到
          四阶矩，建立了更弱但更实用的误差条件。
    \item \textbf{时间分段神经网络架构}：
          \citep{tang2024diffusion}、\citep{tang2025conditional}
          构造了随时间演化的神经网络架构，
          有效平衡了近时间端与远时间端的逼近精度。
    \item \textbf{$s$-conductance profile}：
          \citep{tang2024mala} 提出的马尔可夫链分析新技术，
          提供了介于经典 conductance 与 conductance profile
          之间的精细分析工具。
    \item \textbf{黎曼流形 Bernstein-von Mises 定理}：
          \citep{tang2023robust} 将经典 BvM 推广到黎曼子流形 setting，
          为流形贝叶斯推断提供了严格的理论基础。
    \item \textbf{Girsanov 定理桥梁}：
          \citep{tang2024diffusion}、\citep{tang2025conditional}
          通过 Girsanov 定理建立分数估计误差
          与分布估计误差之间的联系，
          是连接 SDE 理论与统计学习的核心工具。
\end{enumerate}

% ----------------------------------------------------------------
% Standalone guard 结束
\ifdefined\STANDALONE
\else
\end{document}
\fi
% ----------------------------------------------------------------
